%\section{Systematic Uncertainties} 

Systematic uncertainties on \xsecproton have been assessed for detector energy response (including hadron propagation), GENIE modeling (including FSI modeling and modeling of 2p2h effects) and the NuMI flux which is obtained from comparisons of the hadron production model with data from hadron production experiment NA49~\cite{Aliaga:2016oaz}.
% Other sources of uncertainty include an extra $10\%$ on the fraction of DIS and resonant to account for the difference in background predictions from the background constraint techniques,
A second background constraint technique was used in which the DIS and resonant background components are floated in the fit. This result differs from the main background constraint by $10\%$, which we take as a systematic uncertainty. Also, we assess the uncertainties from the efficiency of the Michel electron cut, and the number of nucleons in the nuclear targets.
  Most uncertainties are evaluated by randomly varying the associated parameters
in the simulation within uncertainties and re-extracting
\xsecproton. 
Uncertainties in the beam flux affect the normalization of \xsecproton
and are correlated across \QSqproton bins. The uncertainties on the
signal neutrino interaction and FSI models affect \xsecproton primarily through the efficiency
correction, and are dominated by uncertainties on the resonance
production axial form factor, pion absorption, and pion inelastic scattering. The uncertainties
associated with hadron propagation in the \minerva detector
are evaluated by reweighting signal and background simulation by 10$\%$, $15\%$, $20\%$ for C, Fe and Pb respectively, based on comparisons between Geant4 and measurements of $\pi,p,n$-nucleus cross-section on nuclei ~\cite{doi:10.1146/annurev.nucl.52.050102.090713,PhysRevC.23.2173,ALLARDYCE19731,PhysRevC.53.1745}. The assigned systematic uncertainties are shown in the bottom right of Fig \ref{fig:cross_section_model}.
